\documentclass{article}

\usepackage[spanish]{babel}
\usepackage{amsmath,amssymb}

\title{Conceptos - Análisis de Algoritmos y Teoría de Autómatas}
\author{Benjamín Rivas Beltrán}
\date{\today}

\begin{document}

\maketitle

\section{Conceptos básicos} 

\begin{itemize}
  \item \textbf{Algoritmo:} Procedimiento bien definido, discreto y finito que permite resolver un problema o realizar una tarea específica.
  \item \textbf{Problema:} Especificación de una tarea que se desea realizar, generalmente expresada en términos de entrada y salida.
\end{itemize}

\section{Propiedades de los algoritmos}

\begin{itemize}
  \item \textbf{Correctitud:} Un algoritmo es correcto si cumple las siguientes condiciones:
  \begin{itemize}
    \item \textbf{Resuelve el problema:} Produce la salida correcta para todas las entradas válidas.
    \item \textbf{Es Robusto:} Maneja adecuadamente entradas inválidas o excepcionales.
    \item \textbf{Termina en tiempo finito:} No entra en bucles infinitos y finaliza después de un número finito de pasos
  \end{itemize}

  \item \textbf{Eficiencia:} Un algoritmo es eficiente si utiliza una cantidad razonable de recursos (tiempo y espacio) en relación con el tamaño de la entrada.
  \item \textbf{Complejidad:} La complejidad de un algoritmo se refiere a la cantidad de recursos que consume en función del tamaño de la entrada. Se clasifica en:
  \begin{itemize}
    \item \textbf{Complejidad temporal:} Mide el tiempo de ejecución del algoritmo en función del tamaño de la entrada.
    \item \textbf{Complejidad espacial:} Mide la cantidad de memoria utilizada por el algoritmo en función del tamaño de la entrada.
  \end{itemize}
\end{itemize}

\end{document}